\begin{proof}
Replace every matrix inequality $Y\succeq0$ in Theorem~\ref{I:U:thm:exact} by an equation $Y=UU^*$ with a fresh square matrix $U$ of the appropriate size; Hermitian variables can be represented by their independent real coordinates. The condition $P\succeq0$ is already implied by $P=WW^*$. For complex data, separate real and imaginary parts. All resulting scalar polynomial equations and inequalities have degree at most three. The cubic terms are $B_aPB_a^*$, $B_a^*QB_a$, and $W^*QW$; all PSD factorizations themselves are quadratic. The number of matrices of order $m$ is constant, yielding the stated variable count. Expanding the sums and products produces polynomially many terms in $d,n,k$ and in the binary size of rational input coefficients.

By Theorem~\ref{I:U:thm:exact}, the projection of this feasible set onto $t$ is exactly the closed ray $[\delta_k(H_f),\infty)$. Quantifier elimination over the real algebraic numbers gives a univariate semialgebraic description of this ray, whose finite endpoint is algebraic. Applying sample-point extraction at that endpoint produces an algebraic feasible tuple and hence an optimal algebraic realization. These are standard effective operations over real closed fields \cite{basu2006,renegar1992}.

For rational input the threshold formula is an existential real formula of polynomial size, which proves membership in $\exists\mathbb R$. The stipulated source stability can itself be checked in polynomial bit time: solve $Y-\sum_aA_aYA_a^*=\UId_n$ over the rationals and test $Y\succ0$. A singular linear system is rejected; a positive-definite solution implies stability by conjugating the positive map with $Y^{1/2}$, and stability gives a unique such solution. Thus invalid source encodings can be rejected before the reduction. If $k$ is supplied in binary, the case $k\ge n$ is trivial for $t\ge0$, because the source itself is admissible. Otherwise $k<n$, so the construction is polynomial in the input size. Negative thresholds are rejected before constructing the formula. The case $k=0$ has the sole candidate $G=0$ and is handled by the same Gramian norm test without candidate variables.
\end{proof}
